% -----------------------------------------------------------------------------------------------------
\section{次领头幂次与端点因子化}\label{sec:front-nlp}

在尺度比的次领头幂次（NLP），SCET 因子化公式含有单尺度函数的卷积 $\int_0^1\dd z\,H(z)\langle O(z)\rangle$，一旦各函数被重整化或在 $d=4$ 附近展开，卷积就在端点 $z\to0,1$ 发散。在 SCET$_{\rm I}$ 型问题中，卷积的 $1/\varepsilon^n$ 极点不能由各函数的重整化消去——“an unexpected failure of dimensional regularization and the $\overline{\rm MS}$ scheme”\cite{Liu:2019oav}；在 SCET$_{\rm II}$ 型问题中还需要快度正规子。这一现象最早见于遍举 $B$ 衰变，在 NLP 是普遍的（阈值 Drell--Yan、Higgs 产生与 DIS 的非对角道，含幂次压低部分子道的事件形状，轻夸克圈诱导的 Higgs 振幅）。核心陈述是\cref{prop:thr-nlp}（[P]）与\cref{prop:power-nlp-endpoint}（[P]）；“端点发散在 NLP 普遍存在、各函数的 $\MSbar$ 重整化不足以消除它们”由显式例子确立（第二十一章）。误解 M22：领头幂次的因子化不能直接推广到次领头幂次。

\begin{assumption}[SCET 的模式完备性与 Glauber 模式的缺席 \tagA]\label{asm:front-scet-modes}
(SM1) 对所考虑的过程，所列的硬、（反）硬共线、共线、软与软夸克模式在所述幂次上再现全部贡献的区域（模式集合完备，零箱减除自洽，\cref{def:scet-zero-bin}）。(SM2) Glauber 模式在所述幂次上没有贡献：对衰变到色中性末态与 $e^+e^-$ 湮灭这是自动的；对有入射部分子的过程的部分子系数函数，陈述只涉及部分子截面（强子层次另需该过程的领头幂次因子化定理，如\cref{thm:dy-factorization}）。

地位：[A] 级——SCET 的标准假设，为所涉过程的区域分析所覆盖；\cref{prop:front-hidden-regions}表明完备性必须逐个过程检验。文献：\cite{Liu:2019oav,Liu:2020wbn,Beneke:2022obx,RothsteinStewart2016}。被使用于：\cref{thm:front-hgg,prop:front-ggh,prop:front-nlp-ll,prop:front-nlp-dis,prop:front-sidis-nlp}。解除途径：对每个过程在所述幂次上作区域分析（钳夹分析与幂次计数，第四章、第九章），显式处理 Glauber 区。
\uses{def:scet-modes, def:scet-zero-bin}
\end{assumption}

\paragraph{地位与缺口。}[A] 级：所列过程的区域分析支持它，但没有作为精确陈述的证明；它是本节各 [A] 陈述的瓶颈（误差类型 E-S：模式完备性是标度陈述）。

\begin{definition}[端点发散与重因子化条件]\label{def:front-endpoint}
对单尺度函数的卷积 $\int_0^1\dd z\,H(z)\langle O(z)\rangle$，在重整化或绕 $d=4$ 展开之后于 $z\to0$ 或 $z\to1$ 出现的发散称为\emph{端点发散}；$[\![f(z)]\!]$ 表示 $f$ 在 $z\to0$ 时的领头项（对数除外）。\emph{重因子化（端点因子化）条件}是因子化公式不同项的端点渐近行为之间的全阶恒等式，借助它们可把各项重排为分别无端点发散的项（如\cref{thm:front-hgg}(b) 的条件）。
\uses{def:scet-modes}
\end{definition}

\begin{theorem}[$h\to\gamma\gamma$（$b$ 夸克圈）的重整化端点因子化定理 \tagA]\label{thm:front-hgg}
令 $\lambda=m_b/M_h$；$b$ 夸克贡献 $\mathcal M_b$ 是 $\order{\lambda}$ 的（一次质量插入），两个光子共线扇区之间交换软夸克。SCET 约定（第十五章）：$n_1=(1,0,0,1)$，$n_2=\bar n_1$，$\bar n_1\cdot k_1=\bar n_2\cdot k_2=M_h$；$O_1=\frac{m_b}{e_b^2}\,h\,\mathcal A^\perp_{n_1\mu}\mathcal A^{\perp\mu}_{n_2}$，$O_2(z)=h\,\bar{\mathcal X}_{n_1}\gamma_\perp^\mu\frac{\bar{\sla n}_1}{2}\delta(z\,\bar n_1\cdot k_1+\ii\bar n_1\cdot\partial)\mathcal X_{n_1}\mathcal A^\perp_{n_2\mu}$，$O_3$ 是 Higgs 流 $h\bar{\mathcal X}_{n_1}\mathcal X_{n_2}$ 与两个次领头幂次拉氏量（硬共线场与软夸克 $q_s$ 的耦合）的编时乘积，其矩阵元因子化为两个辐射喷注函数 $J$ 与软夸克函数 $S$ 的卷积，$\langle O_3\rangle=\frac{g_\perp^{\mu\nu}}2\int_0^\infty\frac{\dd\ell_+}{\ell_+}\int_0^\infty\frac{\dd\ell_-}{\ell_-}[J(M_h\ell_+)J(-M_h\ell_-)+J(-M_h\ell_+)J(M_h\ell_-)]S(\ell_+\ell_-)$；$H_2(z)=\bar H_2(z)/[z(1-z)]$。
(a)（裸定理）$\mathcal M_b=H_1\langle O_1\rangle+2\int_0^1\dd z\,H_2(z)\langle O_2(z)\rangle+H_3\langle O_3\rangle$；卷积用解析（快度）正规子 $\eta$ 定义（先对 $\eta$、后对 $\varepsilon$ 展开）\cite{Liu:2019oav}。
(b)（重因子化条件，在 SCET 中以 $D$ 维全阶证明，第二式经软函数的不连续性\cite[第~3~节]{Liu:2020wbn}）
\begin{equation}\label{eq:front-refact}
[\![\bar H_2(z)]\!]=-H_3\,J(zM_h^2),\qquad
[\![\langle O_2(z)\rangle]\!]=-\frac{g_\perp^{\mu\nu}}2\int_0^\infty\frac{\dd\ell_+}{\ell_+}\,J(-M_h\ell_+)\,S(zM_h\ell_+);
\end{equation}
它们蕴含快度正规子依赖在全阶相消。
(c)（重整化定理\cite[式~(4.1)]{Liu:2020wbn}；$\bar z=1-z$）
\begin{equation}\label{eq:front-hgg-ren}
\begin{split}
\mathcal M_b&=H_1(\mu)\langle O_1(\mu)\rangle+2\int_0^1\dd z\Bigl[H_2(z,\mu)\langle O_2(z,\mu)\rangle-\frac{[\![\bar H_2(z,\mu)]\!]}{z}[\![\langle O_2(z,\mu)\rangle]\!]\\
&\hspace{9.5em}-\frac{[\![\bar H_2(\bar z,\mu)]\!]}{\bar z}[\![\langle O_2(\bar z,\mu)\rangle]\!]\Bigr]\\
&\quad+g_\perp^{\mu\nu}\lim_{\sigma\to-1}H_3(\mu)\int_0^{M_h}\frac{\dd\ell_-}{\ell_-}\int_0^{\sigma M_h}\frac{\dd\ell_+}{\ell_+}\,J(M_h\ell_-,\mu)\,J(-M_h\ell_+,\mu)\,S(\ell_+\ell_-,\mu)\Bigr|_{\rm LP}
\end{split}
\end{equation}
在全阶成立；端点截断与重整化不可交换的部分被吸收进 $H_1(\mu)$；各成分满足重整化群方程（多数反常维度已知到两圈）；定理预言三圈对数 $\order{\alpha_b\alpha_s^2L^k}$，$k=6,5,4,3$，$L=\ln(-M_h^2/m_b^2)$。（本定理只涉及部分子振幅。）
% VF-FRO-9
% errtype: E-S
\uses{asm:front-scet-modes, def:front-endpoint, prop:scet-scet2}
\end{theorem}

\begin{proofsketch}
(1)~构造 $\order{\lambda^3}$ 的 SCET 算符基 $O_1,O_2,O_3$ 并与 QCD 匹配（在\cref{asm:front-scet-modes}之下，硬、共线、软夸克区域穷尽领头贡献，见\cref{fig:front-endpoint}(a)）；(2)~求出 $H_2$ 与 $\langle O_2\rangle$ 在 $z\to0$ 的端点行为；(3)~在 $D$ 维证明条件\cref{eq:front-refact}；(4)~用它们重排因子化公式：第二项减去端点渐近，第三项加上截断 $M_h$、$\sigma M_h$（截断的必要性来自共线反常），重复减去的“无穷远箱”并入 $H_1$；(5)~逐项重整化并核对重整化群相容性。完整论证见\cite{Liu:2019oav,Liu:2020wbn}。
\end{proofsketch}

\paragraph{地位与缺口。}[A]：带重因子化条件的 SCET 全阶论证，条件本身在 SCET 内证明；最弱的一步是 SCET 的假设\cref{asm:front-scet-modes}（色中性末态，没有 Glauber 问题）。在\cite{Liu:2019oav}中这些条件只在 NLO 核对、对全阶是猜想。所预言的三圈对数与一个数值计算一致；据\cite{Liu:2020wbn}，对次领头对数（三圈的 $k=5$ 项）它与早先常规方法的预言不一致——这是作者报告的分歧，本书不作裁决。

%%FIGBLOCK fig:front-endpoint%%

\begin{example}[领头双对数与 ${}_2F_2$]\label{ex:front-hgg-ldl}
$h\to\gamma\gamma$ 的 $b$ 圈振幅的领头双对数为\cite{Kotsky:1997rq,Akhoury:2001mz}（三圈见\cite{Penin:2014msa,Liu:2017vkm,Liu:2018czl}；按\cite[式~(3)]{Liu:2019oav}的引述）
\begin{equation}\label{eq:front-ldl}
\mathcal M_b\big|_{\rm LDL}=\mathcal M_0\hat L^2\sum_{n\ge0}\frac{n!}{(2n+2)!}\Bigl(-\frac{C_F\alpha_s}{2\pi}\hat L^2\Bigr)^n
=\mathcal M_0\frac{\hat L^2}{2}\,{}_2F_2\Bigl(1,1;\frac32,2;-\frac{C_F\alpha_s}{8\pi}\hat L^2\Bigr),
\end{equation}
$\hat L=\ln(M_h^2/m_b^2)$，$\mathcal M_0$ 是领头阶的归一化（正比于 $N_c\alpha_by_bm_b$ 与光子极化）。由 $(\frac32)_k=\frac{(2k+1)!}{4^kk!}$ 得 $\frac12\frac{k!}{(\frac32)_k(k+1)!}4^{-k}=\frac{k!}{2(k+1)(2k+1)!}=\frac{k!}{(2k+2)!}$，而 $\frac{(1)_k(1)_k}{(\frac32)_k(2)_k\,k!}=\frac{k!}{(\frac32)_k(k+1)!}$，逐项比较即得第二个等号。与 NLO 比较：$n=0$ 项为 $\frac{\hat L^2}2$，$n=1$ 项为 $-\frac{C_F\alpha_s}{48\pi}\hat L^4=\frac{C_F\alpha_s}{4\pi}\bigl(-\frac{\hat L^4}{12}\bigr)$，正是精确 NLO 振幅\cite[式~(1)]{Liu:2019oav}的领头对数项。\cref{thm:front-hgg}再现了这一级数。
% check: writing/work/W27/s_series.py (B1)-(B4); frontier check OTHERS s4
\uses{thm:front-hgg}
\end{example}

\begin{proposition}[轻夸克诱导的 $gg\to h$ 与重因子化减除方案 \tagA]\label{prop:front-ggh}
同样的结构在带外部色荷的轻夸克圈 $gg\to h$ 中成立；“重因子化减除”（RBS）方案给出没有端点发散的因子化定理，预言三圈对数 $\alpha_s^3\ln^k$（$k=6,\dots,3$），并重求和前三个对数塔 $\alpha_s^n\ln^{2n-k}(-M_h^2/m_b^2)$，$k=0,1,2$（\cite{Liu:2022ajh}，据该文摘要）。常规 QCD 中质量压低振幅的全阶领头对数见\cite{Liu:2017vkm,Liu:2018czl}。坐标空间 SCET 对各成分的分析（均据摘要）：软夸克函数的解析性、重整化与演化\cite{Bodwin:2021cpx}，辐射喷注函数的重整化\cite{Bodwin:2021epw}与规范不变性\cite{Bodwin:2023asf}，共线函数中“隐藏的自能贡献”\cite{Bodwin:2025zzi}（预印本）。$gg\to h$ 软夸克算符的无快度发散定义，以及（借共形对称性）$\gamma\gamma\to h$ 软核与 $B$ 介子光锥分布振幅核之间的全阶关系、并自举到两圈，见\cite{Beneke:2024cpq}（据摘要）。（本命题只涉及部分子振幅。）
% errtype: E-S
\uses{thm:front-hgg, asm:front-scet-modes}
\end{proposition}

\begin{proofsketch}
把\cref{thm:front-hgg}的算符基与重因子化条件推广到带色的外态，并以减除项实现端点重排（RBS 方案）。完整论证见\cite{Liu:2022ajh}。
\end{proofsketch}

\paragraph{地位与缺口。}[A]（只读到摘要层次）；瓶颈同\cref{thm:front-hgg}，即\cref{asm:front-scet-modes}。

\begin{proposition}[次领头幂次阈值 Drell--Yan 的 $d$ 维因子化 \tagP]\label{prop:front-nlp-threshold}
在移去维数正规子之前（$d$ 维），$q\bar q\to\gamma^*+X$ 在 $z=Q^2/\hat s\to1$ 附近对任意次领头幂次有一般的因子化公式，NLP 共线函数由算符匹配定义，并与区域方法核对到 NNLO；若先把共线与软函数在 $d=4$ 附近展开，则“factorization holds only before the dimensional regulator is removed, due to a divergent convolution”，这是超出领头对数的 NLP 重求和的障碍（\cite{Beneke:2019oqx}，据该文摘要）。（本命题只涉及部分子截面。）
\uses{def:thr-limit, thm:thr-refactorization, def:front-endpoint}
\end{proposition}

\paragraph{地位与缺口。}[P]（裸的 $d$ 维形式）；细化了\cref{prop:thr-nlp}（[P]）。

\begin{proposition}[次领头幂次的领头对数重求和与端点因子化条件 \tagA]\label{prop:front-nlp-ll}
在领头对数精度上：(a) Drell--Yan\cite{Beneke:2018gvs}与 $gg\to H$\cite{Beneke:2019mua}的 NLP 阈值重求和（据摘要）；(b) 大 $x$ 处非对角 DIS 由 $d$ 维重因子化得到：区域之间 $1/\varepsilon$ 极点的相消自举出领头对数级数，重现 Vogt 关于大 $x$ 非对角 DGLAP 核的 Bernoulli 数级数，非对角劈裂函数与软夸克 Sudakov 对数是双尺度对象（\cite{Beneke:2020ibj}，据摘要）；(c) “胶子推力”（推力中幂次压低的 $q\bar q$ 对胶子喷注道）：算符层次的端点因子化条件——B 型（反共线）项的 $r\to0$ 极限在全阶等于 A 型（软）项的 $\omega\to\infty$ 极限，作为“prerequisite for endpoint factorization”，在 $\order{\alpha_s}$ 核对；借此得到各项分别无端点发散、由重整化的四维函数表达的因子化公式，以及对端点发散 SCET$_{\rm I}$ 观测量的重整化群领头对数重求和（作者称这是此类观测量的第一个），与早先的 $d$ 维结果一致\cite[第~1--2、4~节]{Beneke:2022obx}；(d) Drell--Yan $qg$ 道 NLP 软函数的重整化与重整化群方程（\cite{Beneke:2025ufd}，据摘要）。相关工作（据摘要）：NLP 阈值对数的次程函／推广 Low--Burnett--Kroll 因子化\cite{Bonocore:2015esa,Bonocore:2016awd}（N$^3$LO 核对见\cite{Bahjat-Abbas:2018hpv}），事件形状的次领头幂次重求和（$H\to gg$ 推力，领头对数\cite{Moult:2018jjd}）与辐射函数\cite{Moult:2019mog}。（本命题只涉及部分子截面与系数函数。）
% errtype: E-S
\uses{asm:front-scet-modes, def:front-endpoint, thm:thr-refactorization}
\end{proposition}

\begin{proofsketch}
在 SCET 中写出 NLP 的算符与区域，导出（(c) 中在算符层次证明）端点渐近之间的全阶关系，据此重排因子化公式，再用单圈反常维度在领头对数精度上求解重整化群方程。完整论证见所引文献。
\end{proofsketch}

\paragraph{地位与缺口。}[A]，只在领头对数精度上（全阶端点条件是论证，(c) 中在单圈核对）；瓶颈\cref{asm:front-scet-modes}；超出领头对数见\cref{conj:front-nlp-beyond-ll}。

\begin{proposition}[次领头幂次 DIS 与 PDF 因子化：阈值 PDF 的双标度结构 \tagA]\label{prop:front-nlp-dis}
在领头对数精度上（预印本），非对角 DIS 道在 $1-x$ 的 NLP（领头扭度）有坐标空间 SCET 因子化：反硬共线 B1 流与软共线（次领头拉氏量）贡献的端点发散由端点因子化关系消解；阈值附近的 PDF 在 NLP 成为双标度对象；在领头幂次，阈值 PDF 与 $\MSbar$ PDF 的 $\xi\to1$ 极限一致，$\MSbar$ 重整化与端点极限可交换；领头对数重求和在 $d$ 维 $\MSbar$ 与重因子化的“端点方案”中都已完成，两种 PDF 定义的关系在双对数精度上全阶给出（\cite{Beneke:2026miy}，预印本；读至摘要与结论）。（本命题只涉及部分子系数函数；强子层次见\cref{thm:pdf-dis-xspace}。）
% errtype: E-S
\uses{asm:front-scet-modes, def:front-endpoint, def:pdf-quark, thm:pdf-dis-xspace}
\end{proposition}

\begin{proofsketch}
在坐标空间 SCET 中匹配 NLP 的 B1 流与次领头拉氏量的贡献，导出其端点渐近之间的关系并重排，再分别在两种方案中重求和领头对数。完整论证见\cite{Beneke:2026miy}（预印本）。
\end{proofsketch}

\paragraph{地位与缺口。}[A]，只在领头对数精度上，且证明仅见于预印本；瓶颈\cref{asm:front-scet-modes}。

\begin{conjecture}[超出领头对数的次领头幂次端点因子化 \tagC]\label{conj:front-nlp-beyond-ll}
对有入射强子、带端点发散的截面（阈值 Drell--Yan 与 DIS 非对角道、含幂次压低道的事件形状），带端点重因子化的 NLP 因子化在超出领头对数的精度上以重整化的四维函数成立，且重因子化条件来自一个一般定理。
\uses{def:front-endpoint, prop:front-nlp-ll}
\end{conjecture}

\paragraph{地位与缺口。}[C]：开放问题，没有超出领头对数的一般论证。
